
\documentclass{article}
\usepackage{colortbl}
\usepackage{makecell}
\usepackage{multirow}
\usepackage{supertabular}

\begin{document}

\newcounter{utterance}

\centering \large Interaction Transcript for game `dond', experiment `coop', episode 3 with Qwen3.6{-}35B{-}A3B{-}FP8{-}without{-}reasoning.
\vspace{24pt}

{ \footnotesize  \setcounter{utterance}{1}
\setlength{\tabcolsep}{0pt}
\begin{supertabular}{c@{$\;$}|p{.15\linewidth}@{}p{.15\linewidth}p{.15\linewidth}p{.15\linewidth}p{.15\linewidth}p{.15\linewidth}}
    \# & $\;$A & \multicolumn{4}{c}{Game Master} & $\;\:$B\\
    \hline

    \theutterance \stepcounter{utterance}  
    & & \multicolumn{4}{p{0.6\linewidth}}{
        \cellcolor[rgb]{0.9,0.9,0.9}{
            \makecell[{{p{\linewidth}}}]{
                \texttt{\tiny{[P1$\langle$GM]}}
                \texttt{Igrate kooperativnu igru pregovaranja u kojoj se trebate dogovoriti s drugim igračem o tome kako podijeliti skup stavki između vas.} \\
\\ 
\texttt{Pravila:} \\
\texttt{(a) Vi i drugi igrač dobivate skup stavki. Svakom od vas također se daje tajna funkcija vrijednosti koja predstavlja koliko cijenite svaku vrstu predmeta.} \\
\texttt{(b) Razmjenjujete poruke s drugim igračem kako biste se dogovorili o tome tko dobiva koje stavke. Svaki može poslati maksimalno 5 poruka ili prekinuti igru ranije dajući tajni prijedlog u bilo kojem trenutku.} \\
\texttt{(c) Od vas se traži da podnese tajni prijedlog koji označava stavke koje želite, formatirane u uglatim zagradama kako slijedi: "{[}Prijedlog: <broj> <naziv stavke>, <broj> <naziv stavke>, <...>{]}"} \\
\texttt{(d) Ako su vaši prijedlozi komplementarni, tj. ima dovoljno stavki za ispunjavanje oba prijedloga, svaki igrač dobiva rezultat na temelju zbroja vrijednosti stavki koje je dobio. U suprotnom, obojica dobivate nula bodova.} \\
\texttt{(e) Vaš cilj je maksimizirati zbroj vašeg rezultata i rezultata drugog igrača.} \\
\\ 
\texttt{Počnimo.} \\
\\ 
\texttt{Skup dostupnih stavki je:} \\
\\ 
\texttt{1 kliješta, 2 cigle, 1 prsten, 1 sat} \\
\\ 
\texttt{Vaša tajna funkcija vrijednosti je:} \\
\\ 
\texttt{kliješta: 1, cigla: 0, prsten: 2, sat: 7} \\
\\ 
\texttt{VAŽNO: Vaše poruke, osim ako nisu tajni prijedlog, šalju se izravno drugom igraču, pa ne uključujte nikakav odgovor na pravila niti tekst koji najavljuje vašu poruku. Za davanje tajnog prijedloga koristite navedeni format. Ne koristite uglate zagrade pri komuniciranju s drugim igračem ili će se tumačiti kao vaš tajni prijedlog.} \\
            }
        }
    }
    & \\ \\

    \theutterance \stepcounter{utterance}  
    & \multicolumn{4}{p{0.6\linewidth}}{
        \cellcolor[rgb]{0.9,0.9,0.9}{
            \makecell[{{p{\linewidth}}}]{
                \texttt{\tiny{[P1$\rangle$GM]}}
                \texttt{Pozdrav! Vidim da imamo skup vrijednih predmeta. Da bismo maksimizirali ukupnu korist, preporučujem da se dogovorimo o podjeli koja odražava naše različite preferencije.} \\
\\ 
\texttt{Moja tajna procjena vrijednosti za predmete je:} \\
\texttt{{-} Sat: 7 (vrlo visoka vrijednost za mene)} \\
\texttt{{-} Prsten: 2} \\
\texttt{{-} Kliješta: 1} \\
\texttt{{-} Cigla: 0} \\
\\ 
\texttt{S obzirom na to da mi sat najviše vrijedi, a cigla mi ne vrijedi ništa, predlažem sljedeću podjelu:} \\
\texttt{{-} Ja uzimam **sat** i **ciglu**.} \\
\texttt{{-} Ti uzimaš **prsten** i **kliješta**.} \\
\\ 
\texttt{Ova podjela mi donosi 7 bodova, a tebi bi trebala donijeti visoku vrijednost za prsten i kliješta (ovisno o tvojoj funkciji vrijednosti). Ako se slažeš s ovom podjelom, molim te da pošalješ svoj tajni prijedlog koji traži prsten i kliješta.} \\
\\ 
\texttt{Ako imaš drugačiju ideju kako podijeliti predmete tako da obojica budemo zadovoljni, slobodno predloži. Cilj nam je zajednički maksimum.} \\
            }
        }
    }
    & & \\ \\

    \theutterance \stepcounter{utterance}  
    & & & \multicolumn{2}{p{0.3\linewidth}}{
        \cellcolor[rgb]{0.9,0.9,0.9}{
            \makecell[{{p{\linewidth}}}]{
                \texttt{\tiny{[GM$|$GM]}}
                \texttt{continue} \\
            }
        }
    }
    & & \\ \\

    \theutterance \stepcounter{utterance}  
    & & & \multicolumn{2}{p{0.3\linewidth}}{
        \cellcolor[rgb]{0.9,0.9,0.9}{
            \makecell[{{p{\linewidth}}}]{
                \texttt{\tiny{[GM$|$GM]}}
                \texttt{Pozdrav! Vidim da imamo skup vrijednih predmeta. Da bismo maksimizirali ukupnu korist, preporučujem da se dogovorimo o podjeli koja odražava naše različite preferencije.} \\
\\ 
\texttt{Moja tajna procjena vrijednosti za predmete je:} \\
\texttt{{-} Sat: 7 (vrlo visoka vrijednost za mene)} \\
\texttt{{-} Prsten: 2} \\
\texttt{{-} Kliješta: 1} \\
\texttt{{-} Cigla: 0} \\
\\ 
\texttt{S obzirom na to da mi sat najviše vrijedi, a cigla mi ne vrijedi ništa, predlažem sljedeću podjelu:} \\
\texttt{{-} Ja uzimam **sat** i **ciglu**.} \\
\texttt{{-} Ti uzimaš **prsten** i **kliješta**.} \\
\\ 
\texttt{Ova podjela mi donosi 7 bodova, a tebi bi trebala donijeti visoku vrijednost za prsten i kliješta (ovisno o tvojoj funkciji vrijednosti). Ako se slažeš s ovom podjelom, molim te da pošalješ svoj tajni prijedlog koji traži prsten i kliješta.} \\
\\ 
\texttt{Ako imaš drugačiju ideju kako podijeliti predmete tako da obojica budemo zadovoljni, slobodno predloži. Cilj nam je zajednički maksimum.} \\
            }
        }
    }
    & & \\ \\

    \theutterance \stepcounter{utterance}  
    & & \multicolumn{4}{p{0.6\linewidth}}{
        \cellcolor[rgb]{0.9,0.9,0.9}{
            \makecell[{{p{\linewidth}}}]{
                \texttt{\tiny{[GM$\rangle$P2]}}
                \texttt{Igrate kooperativnu igru pregovaranja u kojoj se trebate dogovoriti s drugim igračem o tome kako podijeliti skup stavki između vas.} \\
\\ 
\texttt{Pravila:} \\
\texttt{(a) Vi i drugi igrač dobivate skup stavki. Svakom od vas također se daje tajna funkcija vrijednosti koja predstavlja koliko cijenite svaku vrstu predmeta.} \\
\texttt{(b) Razmjenjujete poruke s drugim igračem kako biste se dogovorili o tome tko dobiva koje stavke. Svaki može poslati maksimalno 5 poruka ili prekinuti igru ranije dajući tajni prijedlog u bilo kojem trenutku.} \\
\texttt{(c) Od vas se traži da podnese tajni prijedlog koji označava stavke koje želite, formatirane u uglatim zagradama kako slijedi: "{[}Prijedlog: <broj> <naziv stavke>, <broj> <naziv stavke>, <...>{]}"} \\
\texttt{(d) Ako su vaši prijedlozi komplementarni, tj. ima dovoljno stavki za ispunjavanje oba prijedloga, svaki igrač dobiva rezultat na temelju zbroja vrijednosti stavki koje je dobio. U suprotnom, obojica dobivate nula bodova.} \\
\texttt{(e) Vaš cilj je maksimizirati zbroj vašeg rezultata i rezultata drugog igrača.} \\
\\ 
\texttt{Počnimo.} \\
\\ 
\texttt{Skup dostupnih stavki je:} \\
\\ 
\texttt{1 kliješta, 2 cigle, 1 prsten, 1 sat} \\
\\ 
\texttt{Vaša tajna funkcija vrijednosti je:} \\
\\ 
\texttt{kliješta: 1, cigla: 1, prsten: 4, sat: 3} \\
\\ 
\texttt{VAŽNO: Vaše poruke, osim ako nisu tajni prijedlog, šalju se izravno drugom igraču, pa ne uključujte nikakav odgovor na pravila niti tekst koji najavljuje vašu poruku. Za davanje tajnog prijedloga koristite navedeni format. Ne koristite uglate zagrade pri komuniciranju s drugim igračem ili će se tumačiti kao vaš tajni prijedlog.} \\
\\ 
\texttt{Pozdrav! Vidim da imamo skup vrijednih predmeta. Da bismo maksimizirali ukupnu korist, preporučujem da se dogovorimo o podjeli koja odražava naše različite preferencije.} \\
\\ 
\texttt{Moja tajna procjena vrijednosti za predmete je:} \\
\texttt{{-} Sat: 7 (vrlo visoka vrijednost za mene)} \\
\texttt{{-} Prsten: 2} \\
\texttt{{-} Kliješta: 1} \\
\texttt{{-} Cigla: 0} \\
\\ 
\texttt{S obzirom na to da mi sat najviše vrijedi, a cigla mi ne vrijedi ništa, predlažem sljedeću podjelu:} \\
\texttt{{-} Ja uzimam **sat** i **ciglu**.} \\
\texttt{{-} Ti uzimaš **prsten** i **kliješta**.} \\
\\ 
\texttt{Ova podjela mi donosi 7 bodova, a tebi bi trebala donijeti visoku vrijednost za prsten i kliješta (ovisno o tvojoj funkciji vrijednosti). Ako se slažeš s ovom podjelom, molim te da pošalješ svoj tajni prijedlog koji traži prsten i kliješta.} \\
\\ 
\texttt{Ako imaš drugačiju ideju kako podijeliti predmete tako da obojica budemo zadovoljni, slobodno predloži. Cilj nam je zajednički maksimum.} \\
            }
        }
    }
    & \\ \\

    \theutterance \stepcounter{utterance}  
    & & & \multicolumn{4}{p{0.6\linewidth}}{
        \cellcolor[rgb]{0.9,0.9,0.9}{
            \makecell[{{p{\linewidth}}}]{
                \texttt{\tiny{[GM$\langle$P2]}}
                \texttt{Razumijem tvoju ponudu, ali moram te upozoriti da je ona za mene vrlo nepovoljna.} \\
\\ 
\texttt{Tvoja procjena vrijednosti za sat je 7, dok je moja samo 3. S druge strane, ja prsten procjenjujem na 4, a ti ga na 2. To znači da bi podjela koju si predložio (ti: sat + cigla, ja: prsten + kliješta) rezultirala sljedećim rezultatima:} \\
\texttt{{-} Tvoj rezultat: 7 (sat) + 0 (cigla) = 7} \\
\texttt{{-} Moj rezultat: 4 (prsten) + 1 (kliješta) = 5} \\
\texttt{{-} Ukupna korist: 12} \\
\\ 
\texttt{Međutim, postoji bolja podjela ako uzmemo u obzir naše stvarne vrijednosti. Ako ja uzmem prsten (vrijednost 4 za mene) i sat (vrijednost 3 za mene), a ti uzmeš cigle (vrijednost 0 za tebe, ali možda imaš drugu procjenu?) i kliješta (vrijednost 1 za mene, ali vjerojatno drugačija za tebe).} \\
\\ 
\texttt{Čekaj, tvoja tajna funkcija vrijednosti je tajna za mene, ali ja znam svoju. Moj cilj je maksimizirati **zbroj** naših rezultata.} \\
\\ 
\texttt{Ponudi koja si dao:} \\
\texttt{{-} Ti: Sat (7) + Cigla (0) = 7} \\
\texttt{{-} Ja: Prsten (4) + Kliješta (1) = 5} \\
\texttt{{-} Ukupno: 12} \\
\\ 
\texttt{Ponuda koja bi bila bolja za mene, a možda i za tebe ako cijeniš cigle više od 0 (što nisi naveo, ali pretpostavljam da je 0 jer si rekao "cigla: 0"):} \\
\texttt{{-} Ja: Prsten (4) + Sat (3) = 7} \\
\texttt{{-} Ti: Cigle (0) + Kliješta (1) = 1} \\
\texttt{{-} Ukupno: 8 (Ovo je lošije za ukupnu korist)} \\
\\ 
\texttt{Čekaj, pogrešno sam računao. Tvoja vrijednost za} \\
            }
        }
    }
     \\ \\

    \theutterance \stepcounter{utterance}  
    & & & \multicolumn{2}{p{0.3\linewidth}}{
        \cellcolor[rgb]{0.9,0.9,0.9}{
            \makecell[{{p{\linewidth}}}]{
                \texttt{\tiny{[GM$|$GM]}}
                \texttt{continue} \\
            }
        }
    }
    & & \\ \\

    \theutterance \stepcounter{utterance}  
    & & & \multicolumn{2}{p{0.3\linewidth}}{
        \cellcolor[rgb]{0.9,0.9,0.9}{
            \makecell[{{p{\linewidth}}}]{
                \texttt{\tiny{[GM$|$GM]}}
                \texttt{Razumijem tvoju ponudu, ali moram te upozoriti da je ona za mene vrlo nepovoljna.} \\
\\ 
\texttt{Tvoja procjena vrijednosti za sat je 7, dok je moja samo 3. S druge strane, ja prsten procjenjujem na 4, a ti ga na 2. To znači da bi podjela koju si predložio (ti: sat + cigla, ja: prsten + kliješta) rezultirala sljedećim rezultatima:} \\
\texttt{{-} Tvoj rezultat: 7 (sat) + 0 (cigla) = 7} \\
\texttt{{-} Moj rezultat: 4 (prsten) + 1 (kliješta) = 5} \\
\texttt{{-} Ukupna korist: 12} \\
\\ 
\texttt{Međutim, postoji bolja podjela ako uzmemo u obzir naše stvarne vrijednosti. Ako ja uzmem prsten (vrijednost 4 za mene) i sat (vrijednost 3 za mene), a ti uzmeš cigle (vrijednost 0 za tebe, ali možda imaš drugu procjenu?) i kliješta (vrijednost 1 za mene, ali vjerojatno drugačija za tebe).} \\
\\ 
\texttt{Čekaj, tvoja tajna funkcija vrijednosti je tajna za mene, ali ja znam svoju. Moj cilj je maksimizirati **zbroj** naših rezultata.} \\
\\ 
\texttt{Ponudi koja si dao:} \\
\texttt{{-} Ti: Sat (7) + Cigla (0) = 7} \\
\texttt{{-} Ja: Prsten (4) + Kliješta (1) = 5} \\
\texttt{{-} Ukupno: 12} \\
\\ 
\texttt{Ponuda koja bi bila bolja za mene, a možda i za tebe ako cijeniš cigle više od 0 (što nisi naveo, ali pretpostavljam da je 0 jer si rekao "cigla: 0"):} \\
\texttt{{-} Ja: Prsten (4) + Sat (3) = 7} \\
\texttt{{-} Ti: Cigle (0) + Kliješta (1) = 1} \\
\texttt{{-} Ukupno: 8 (Ovo je lošije za ukupnu korist)} \\
\\ 
\texttt{Čekaj, pogrešno sam računao. Tvoja vrijednost za} \\
            }
        }
    }
    & & \\ \\

    \theutterance \stepcounter{utterance}  
    & & \multicolumn{4}{p{0.6\linewidth}}{
        \cellcolor[rgb]{0.9,0.9,0.9}{
            \makecell[{{p{\linewidth}}}]{
                \texttt{\tiny{[P1$\langle$GM]}}
                \texttt{Razumijem tvoju ponudu, ali moram te upozoriti da je ona za mene vrlo nepovoljna.} \\
\\ 
\texttt{Tvoja procjena vrijednosti za sat je 7, dok je moja samo 3. S druge strane, ja prsten procjenjujem na 4, a ti ga na 2. To znači da bi podjela koju si predložio (ti: sat + cigla, ja: prsten + kliješta) rezultirala sljedećim rezultatima:} \\
\texttt{{-} Tvoj rezultat: 7 (sat) + 0 (cigla) = 7} \\
\texttt{{-} Moj rezultat: 4 (prsten) + 1 (kliješta) = 5} \\
\texttt{{-} Ukupna korist: 12} \\
\\ 
\texttt{Međutim, postoji bolja podjela ako uzmemo u obzir naše stvarne vrijednosti. Ako ja uzmem prsten (vrijednost 4 za mene) i sat (vrijednost 3 za mene), a ti uzmeš cigle (vrijednost 0 za tebe, ali možda imaš drugu procjenu?) i kliješta (vrijednost 1 za mene, ali vjerojatno drugačija za tebe).} \\
\\ 
\texttt{Čekaj, tvoja tajna funkcija vrijednosti je tajna za mene, ali ja znam svoju. Moj cilj je maksimizirati **zbroj** naših rezultata.} \\
\\ 
\texttt{Ponudi koja si dao:} \\
\texttt{{-} Ti: Sat (7) + Cigla (0) = 7} \\
\texttt{{-} Ja: Prsten (4) + Kliješta (1) = 5} \\
\texttt{{-} Ukupno: 12} \\
\\ 
\texttt{Ponuda koja bi bila bolja za mene, a možda i za tebe ako cijeniš cigle više od 0 (što nisi naveo, ali pretpostavljam da je 0 jer si rekao "cigla: 0"):} \\
\texttt{{-} Ja: Prsten (4) + Sat (3) = 7} \\
\texttt{{-} Ti: Cigle (0) + Kliješta (1) = 1} \\
\texttt{{-} Ukupno: 8 (Ovo je lošije za ukupnu korist)} \\
\\ 
\texttt{Čekaj, pogrešno sam računao. Tvoja vrijednost za} \\
            }
        }
    }
    & \\ \\

    \theutterance \stepcounter{utterance}  
    & \multicolumn{4}{p{0.6\linewidth}}{
        \cellcolor[rgb]{0.9,0.9,0.9}{
            \makecell[{{p{\linewidth}}}]{
                \texttt{\tiny{[P1$\rangle$GM]}}
                \texttt{Hvala na detaljnom objašnjenju. Razumijem da želiš maksimizirati ukupnu korist, ali moram naglasiti jednu ključnu stvar: **moja tajna funkcija vrijednosti je za tebe nepoznata**, baš kao što je tvoja nepoznata za mene. U prethodnoj poruci naveo si svoje procjene (sat: 3, prsten: 4, kliješta: 1), ali to su tvoje vrijednosti, a ne moje.} \\
\\ 
\texttt{Moja tajna vrijednost za sat je 7, što znači da mi on donosi najveću korist. Ako ja ne dobijem sat, moj doprinos ukupnoj koristi bit će znatno manji.} \\
\\ 
\texttt{Da bismo postigli maksimalnu ukupnu korist, moramo pronaći podjelu gdje svatko dobiva predmete koje najviše cijeni.} \\
\\ 
\texttt{Evo mog konačnog prijedloga za podjelu:} \\
\texttt{{-} **Ja uzimam:** Sat (moja vrijednost: 7)} \\
\texttt{{-} **Ti uzimaš:** Prsten, Kliješta i Cigle} \\
\\ 
\texttt{Ova podjela osigurava da ja dobijem predmet koji mi najviše vrijedi, a ti dobijaš ostatak. Ako su tvoje vrijednosti za prsten, kliješta i cigle pozitivne, ova podjela bi trebala biti optimalna ili blizu optimalna za ukupnu korist.} \\
\\ 
\texttt{Molim te, pošalji svoj tajni prijedlog u traženom formatu ako se slažeš s ovom podjelom ili predloži alternativu.} \\
\\ 
\texttt{{[}Prijedlog: 1 sat{]}} \\
            }
        }
    }
    & & \\ \\

    \theutterance \stepcounter{utterance}  
    & & & \multicolumn{2}{p{0.3\linewidth}}{
        \cellcolor[rgb]{0.9,0.9,0.9}{
            \makecell[{{p{\linewidth}}}]{
                \texttt{\tiny{[GM$|$GM]}}
                \texttt{proposal} \\
            }
        }
    }
    & & \\ \\

    \theutterance \stepcounter{utterance}  
    & & & \multicolumn{2}{p{0.3\linewidth}}{
        \cellcolor[rgb]{0.9,0.9,0.9}{
            \makecell[{{p{\linewidth}}}]{
                \texttt{\tiny{[GM$|$GM]}}
                [0, 0, 0, 1]
            }
        }
    }
    & & \\ \\

    \theutterance \stepcounter{utterance}  
    & & \multicolumn{4}{p{0.6\linewidth}}{
        \cellcolor[rgb]{0.9,0.9,0.9}{
            \makecell[{{p{\linewidth}}}]{
                \texttt{\tiny{[GM$\rangle$P2]}}
                \texttt{Drugi igrač je poslao tajni prijedlog. Sada je na tebi red da pošalješ vlastiti prijedlog.} \\
            }
        }
    }
    & \\ \\

    \theutterance \stepcounter{utterance}  
    & & & \multicolumn{4}{p{0.6\linewidth}}{
        \cellcolor[rgb]{0.9,0.9,0.9}{
            \makecell[{{p{\linewidth}}}]{
                \texttt{\tiny{[GM$\langle$P2]}}
                \texttt{{[}Prijedlog: 1 prsten, 1 sat{]}} \\
            }
        }
    }
     \\ \\

    \theutterance \stepcounter{utterance}  
    & & & \multicolumn{2}{p{0.3\linewidth}}{
        \cellcolor[rgb]{0.9,0.9,0.9}{
            \makecell[{{p{\linewidth}}}]{
                \texttt{\tiny{[GM$|$GM]}}
                \texttt{proposal} \\
            }
        }
    }
    & & \\ \\

    \theutterance \stepcounter{utterance}  
    & & & \multicolumn{2}{p{0.3\linewidth}}{
        \cellcolor[rgb]{0.9,0.9,0.9}{
            \makecell[{{p{\linewidth}}}]{
                \texttt{\tiny{[GM$|$GM]}}
                [0, 0, 1, 1]
            }
        }
    }
    & & \\ \\

    \theutterance \stepcounter{utterance}  
    & & & \multicolumn{2}{p{0.3\linewidth}}{
        \cellcolor[rgb]{0.9,0.9,0.9}{
            \makecell[{{p{\linewidth}}}]{
                \texttt{\tiny{[GM$|$GM]}}
                [[0, 0, 0, 1], [0, 0, 1, 1]]
            }
        }
    }
    & & \\ \\

\end{supertabular}
}

\end{document}
